\section{システム概要}
\label{sec:system}

\begin{figure}[t]
  \centering
  \resizebox{\linewidth}{!}{% システム概要。すべて本稿のために作成した TikZ 図である。
\begin{tikzpicture}[
  font=\sffamily\scriptsize,
  card/.style={draw=black!35, rounded corners=2.2pt, align=center,
    minimum height=24mm, inner sep=3pt, line width=0.45pt},
  pipe/.style={-{Stealth[length=2.2mm,width=1.45mm]}, line width=0.75pt, draw=black!72},
  tag/.style={rounded corners=1.5pt, inner xsep=3pt, inner ysep=1.2pt,
    font=\sffamily\fontsize{6.4}{7.4}\selectfont},
  note/.style={font=\sffamily\fontsize{6.3}{7.2}\selectfont, text=black!67, align=center},
]
  \node[card, fill=softgray!70, minimum width=21mm] (video) {
    \textbf{単一カメラ動画}\\[0.4em]
    RGB 系列\\
    + 会場 ID
  };
  \node[card, fill=gsblue!10, draw=gsblue, minimum width=25mm,
    right=3.2mm of video] (recon) {
    \textbf{再構成}\\[0.4em]
    SfM カメラ\\
    + 3DGS モデル
  };
  \node[card, fill=courtgreen!10, draw=courtgreen, minimum width=30mm,
    right=3.2mm of recon] (align) {
    \textbf{メートル系}\\[0.1em]
    \textbf{コートアライメント}\\[0.4em]
    fit 証拠で当てはめ\\
    holdout で凍結検証
  };
  \node[card, fill=softgray!70, minimum width=28mm,
    right=3.2mm of align] (render) {
    \textbf{視点生成と}\\[0.1em]
    \textbf{レンダリング}\\[0.4em]
    SfM 制約付き軌道\\
    弧長一様サンプル
  };
  \node[card, fill=lossred!8, draw=lossred, minimum width=28mm,
    right=3.2mm of render] (labels) {
    \textbf{整合教師}\\[0.4em]
    \KPName{} + 領域 + 線\\
    + 線の意味 + 姿勢
  };
  \node[card, fill=poseorange!10, draw=poseorange, minimum width=23mm,
    right=3.2mm of labels] (model) {
    \textbf{コート検出器}\\[0.4em]
    共有特徴\\
    + 密予測と姿勢
  };

  \draw[pipe] (video) -- (recon);
  \draw[pipe, draw=gsblue] (recon) -- (align);
  \draw[pipe, draw=courtgreen] (align) -- (render);
  \draw[pipe] (render) -- (labels);
  \draw[pipe, draw=lossred] (labels) -- (model);

  % 二つの検証境界。ここを通過しないシーンは後段へ流れない。
  \node[boundary, fit=(recon), inner sep=2.4mm] (scene_boundary) {};
  \node[tag, fill=gsblue!12, draw=gsblue, anchor=south west]
    at (scene_boundary.north west) {境界 1: 公開シーン出力};
  \node[boundary, draw=courtgreen, fit=(align), inner sep=2.4mm] (metric_boundary) {};
  \node[tag, fill=courtgreen!12, draw=courtgreen, anchor=south west]
    at (metric_boundary.north west) {境界 2: アライメント受理};

  % 下段: 同一変換が画像と教師の両方を支配する。
  \node[note, below=6.5mm of labels, text width=52mm] (share) {
    一つの変換グラフが\\画像と教師の両方を支配する
  };
  \draw[-{Stealth[length=1.8mm,width=1.2mm]}, draw=courtgreen, dashed, line width=0.6pt]
    (metric_boundary.south) .. controls +(0,-6mm) and +(-14mm,2mm) .. (share.west);
  \draw[-{Stealth[length=1.8mm,width=1.2mm]}, draw=courtgreen, dashed, line width=0.6pt]
    (share.east) .. controls +(14mm,-2mm) and +(0,-6mm) .. (labels.south);

  \node[note, below=14.5mm of video, text width=30mm, anchor=north west]
    at (video.south west) {人手の費用は\\シーン単位の検証に移る};
\end{tikzpicture}
}
  \caption{\textbf{システム概要。}
  単一カメラの動画を再構成し、得られた 3DGS シーンを
  メートル系のコートアライメントへ通す。
  アライメントが受理されたシーンだけが、
  制御された三次元軌道からのレンダリングと
  整合した教師生成を許される。
  人手の検証費用は画像あたりではなくシーンあたりに一度だけ発生し、
  その一つの変換が画像と全教師を同時に支配する。
  }
  \label{fig:system}
\end{figure}

システムは二つの独立した境界を持つ（\cref{fig:system}）。

\paragraph{シーン境界}
再構成との境界は、公開された標準シーン出力だけを受理する。
そこには画像、校正済みカメラ、疎な三次元点、明示的な座標規約、
レンダリング可能なガウスモデルが含まれる。
コート側の処理は再構成の内部実装や隠れたチェックポイント配置を参照せず、
公開された二つのコマンド（シーン出力と任意カメラのレンダリング）のみを介する。
出力の検証ではスキーマ、カメラ識別子、形状、型、
有限性、内部パラメータ、回転の直交性、相互変換の整合性を確認する。

\paragraph{メートル境界}
アライメントは、線証拠、尺度、トポロジー、相互変換、
holdout 残差のすべてが型付きの判定を通ったときだけシーンを受理する。
受理されたシーンは
\(\mathcal{L}=\{\T{\Scene_m}{\Court_i},\T{\Court_i}{\Scene_m},
\mathcal{B},\mathcal{M}\}_{i=1}^{N_c}\)
を公開する。
ここで \(\mathcal{B}\) は複数コートの配置範囲、
\(\mathcal{M}\) は fit と holdout の指標である。
下流のデータ生成はこの権威を再推定・再解釈しない。

\paragraph{生成と学習}
受理後は、軌道候補の列挙、弧長一様サンプリング、
軌道グループ単位の分割、任意カメラのレンダリング、
教師の生成が順に行われる。
学習側は、対象コートを基準にカメラ端の向きを正準化し、
共有特徴から \KPName、領域、線、線の意味、姿勢を予測する。

この構成で重要なのは、\emph{どの段階も欠けた権威を静かに補わない}ことである。
アライメントが失敗したシーンは棄却され、代替の一コート仮定を置かない。
holdout が崩れた変換は再当てはめせず失敗とする。
教師の順序が食い違う sample は学習へ入らずエラーになる。
再構成の外挿領域では画質が保証されないため、
軌道は復元カメラの凸包を内側へオフセットした領域へ解析的に制限する。

\begin{table}[t]
  \centering
  \caption{\textbf{段階ごとの権威と失敗時の挙動。}
  どの段階も、欠けた権威を下流で静かに補わない。
  }
  \label{tab:component_contracts}
  \small
  \begin{tabularx}{\linewidth}{@{}p{0.16\linewidth}p{0.30\linewidth}X@{}}
    \toprule
    段階 & 公開する権威 & 推測せず棄却するもの \\
    \midrule
    再構成 & シーンスキーマ、復元カメラ、内部パラメータ、
      正規化点群、レンダ可能な 3DGS &
      隠れた再構成経路、壊れた配列、非直交回転、公開コマンドの欠落 \\
    アライメント & 共通メートル尺度、相互のシーン--コート変換、
      トポロジー、fit/holdout 指標 &
      コート一面への縮退、holdout の再当てはめ、
      尺度不一致、意味セグメント被覆の不足 \\
    データ生成 & RGB と深度、対象コート、カメラ変換、
      \KPName の物理番号、幾何妥当性と描画支持、軌道グループ分割 &
      スキーマ変換、最近傍による点の同一視、
      フレーム単位の分割漏洩、視点の差し替え \\
    検出器 & 型付き密教師、姿勢ベクトル、補正後の内部パラメータ、
      有効領域マスク &
      未対応 source や未対応補正、単一ピーク前提での多峰キーポイント、
      姿勢枝またはキーポイント枝の欠落 \\
    \bottomrule
  \end{tabularx}
\end{table}

